\documentclass[11pt]{article}
\usepackage[margin=0.9in]{geometry}
\usepackage{amsmath,amssymb,amsthm}
\usepackage{enumitem}
\usepackage{parskip}
 
\setlength{\parindent}{0pt}
\setlength{\parskip}{6pt}
 
\newcommand{\dydx}{\dfrac{dy}{dx}}
\newcommand{\Lap}{\mathcal{L}}
 
\title{MTH 4110 --- Differential Equations Study Notes}
\author{}
\date{}
 
\begin{document}
\maketitle
 
\section*{Separation of Variables}
An ODE of the form $g(y)\,\dydx = h(x)$ has the implicit solution given by $G(y) = H(x) + C$ where $G(y) = \int g(y)\,dy$ and $H(x) = \int h(x)\,dx$.
 
\textit{i.e.}\ for $\dydx = (3y-1)^2 \sin(2x)$:
\[
\int (3y-1)^{-2}\,dy = \int \sin(2x)\,dx \;\;\Longrightarrow\;\; \frac{-1}{3(3y-1)} = -\tfrac{1}{2}\cos(2x) + C.
\]
 
\section*{Integrating Factor (Linear First-Order)}
Consider the first-order linear ODE of the form $\dydx + p(x)y = f(x)$. We define the integrating factor as $\mu(x) = e^{\int p(x)\,dx}$ and then multiply through by $\mu(x)$.
 
\textit{i.e.}\ take $x\dydx - 2y = x^3 e^{3x}$, with $y(1) = e^3$. Rewrite as $\dydx - \tfrac{2}{x}y = x^2 e^{3x}$. Integrating factor: $\mu(x) = e^{\int -\frac{2}{x}\,dx} = e^{-2\ln x} = x^{-2}$. So:
\[
x^{-2}\dydx - 2x^{-3}y = e^{3x} \;\;\Longrightarrow\;\; \tfrac{d}{dx}\!\left[x^{-2}y\right] = \tfrac{d}{dx}\!\left[\tfrac{1}{3}e^{3x}\right] \;\;\Longrightarrow\;\; y = \frac{x^2 e^{3x}}{3} + Cx^2.
\]
Using $y(1) = e^3$: $\,C = \tfrac{2}{3}e^3$.
 
\section*{Integrating Factor for Nonexact Equations}
Given a first-order ODE of the form $M(x,y) + N(x,y)y' = 0$, it is \emph{exact} if there exists $\psi(x,y)$ such that $\psi_x = M$ and $\psi_y = N$. The previous ODE is exact if and only if $M_y = N_x$ for all $x,y \in \mathbb{R}$. If $M_y = N_x$, then $\psi(x,y) = \int M\,dx = F(x,y) + h(y)$. Take $\psi_y = \tfrac{d}{dy}\int M\,dx \stackrel{\text{set}}{=} N$. After solving for $h'(y)$, plug back into $\psi(x,y)$. Then the ODE becomes $\tfrac{d}{dx}[\psi(x,y)] = 0 \Rightarrow \psi(x,y) = C$.
 
If $M_y \neq N_x$, use the formula below and then use the result as the integrating factor:
\begin{itemize}[leftmargin=*]
    \item If $\dfrac{M_y - N_x}{N}$ is a function of only $x$, then $\mu(x) = e^{\int \frac{M_y - N_x}{N}\,dx}$.
    \item If $\dfrac{N_x - M_y}{M}$ is a function of only $y$, then $\mu(y) = e^{\int \frac{N_x - M_y}{M}\,dy}$.
\end{itemize}
 
\textit{i.e.}\ for $6x^2 y + (4x^3 + 15y^3)y' = 0$ with $y(2) = 1$: here $M = 6x^2 y$, $N = 4x^3 + 15y^3$, so $M_y = 6x^2$, $N_x = 12x^2$. Since $\frac{N_x - M_y}{M} = \frac{6x^2}{6x^2 y} = \frac{1}{y}$ depends only on $y$, the integrating factor is $\mu(y) = e^{\int \frac{1}{y}\,dy} = y$. Multiplying gives an exact equation; solving yields $2x^3 y^2 + 3y^5 = 19$.
 
\section*{Characteristic Equation (Constant Coefficients)}
If we have a linear homogeneous ODE with constant coefficients, we can look at the roots of the characteristic polynomial to determine the general solution. For $ay'' + by' + cy = 0$, the characteristic equation is $ar^2 + br + c = 0$.
\begin{itemize}[leftmargin=*]
    \item Two distinct real roots: $y(x) = c_1 e^{r_1 x} + c_2 e^{r_2 x}$.
    \item Repeated root: $y(x) = c_1 e^{rx} + c_2 x e^{rx}$.
    \item Complex roots $\alpha \pm \beta i$: $y(x) = e^{\alpha x}\bigl(c_1 \cos(\beta x) + c_2 \sin(\beta x)\bigr)$.
\end{itemize}
 
\textit{i.e.}\ for $y'' + 4y' - 12y = 0$ with $y(0) = 2$, $y'(0) = 12$: characteristic equation $r^2 + 4r - 12 = (r+6)(r-2) = 0$, so $r = -6, 2$. General solution: $y(x) = c_1 e^{-6x} + c_2 e^{2x}$. Applying ICs gives $c_1 = -1$, $c_2 = 3$, so $y(x) = -e^{-6x} + 3e^{2x}$.
 
For higher-order ODEs, each root contributes according to its multiplicity. \textit{i.e.}\ if the characteristic equation factors as $(r+1)^3(r^2 - 4r + 5)(r^2 + 4)^2 = 0$, the roots are $r = -1$ (mult.\ 3), $r = 2 \pm i$ (each mult.\ 1), $r = \pm 2i$ (each mult.\ 2), giving order 9 and
\[
y = c_1 e^{-x} + c_2 x e^{-x} + c_3 x^2 e^{-x} + c_4 e^{2x}\sin x + c_5 e^{2x}\cos x + c_6 \sin(2x) + c_7 x\sin(2x) + c_8 \cos(2x) + c_9 x\cos(2x).
\]
 
\section*{Wronskian and Fundamental Sets}
For $y_1, y_2$ on an interval $I$:
\[
W[y_1, y_2](x) = \begin{vmatrix} y_1 & y_2 \\ y_1' & y_2' \end{vmatrix} = y_1 y_2' - y_2 y_1'.
\]
If $W \neq 0$ on $I$, then $\{y_1, y_2\}$ forms a fundamental set of solutions. Generalizes to $n$ functions via an $n \times n$ determinant.
 
\textit{i.e.}\ for three degree-1 polynomials $y_1 = Ax + B$, $y_2 = Cx + D$, $y_3 = Ex + F$, the third row of $W[y_1,y_2,y_3]$ is all zeros (since each $y_i'' = 0$), so $W = 0$ identically. Thus three linear functions cannot form a fundamental set for a third-order linear ODE.
 
\section*{Reduction of Order}
Given one solution $y_1$ to $y'' + p(x)y' + q(x)y = 0$, look for a second solution of the form $y_2 = v(x)y_1(x)$. Letting $w = v'$, $w$ satisfies the first-order linear ODE
\[
y_1 w' + (2y_1' + p y_1)w = 0.
\]
Solve for $w$, then $v = \int w\,dx$, and finally $y_2 = v y_1$.
 
\textit{i.e.}\ for $x^2 y'' - 7xy' + 16y = 0$ on $x > 0$ with $y_1 = x^4$: standard form gives $p = -7/x$. Then
\[
x^4 w' + \bigl(2(4x^3) - \tfrac{7}{x}(x^4)\bigr)w = 0 \;\Rightarrow\; x^4 w' + x^3 w = 0 \;\Rightarrow\; w = \tfrac{A}{x}.
\]
So $v = A\ln x + B$. Choosing $A = 1$, $B = 0$ gives $y_2 = (\ln x)x^4$.
 
\section*{Undetermined Coefficients}
For $ay'' + by' + cy = g(x)$ where $g$ is a polynomial, exponential, sine/cosine, or product thereof, guess $y_p$ of the same form with unknown coefficients. If any term in the guess overlaps with the complementary solution $y_c$, multiply by $x^s$ where $s$ is the smallest positive integer eliminating the overlap.
 
\textit{i.e.}\ for $y'' + y' - 2y = xe^{4x}$: $r^2 + r - 2 = (r+2)(r-1) = 0$ gives $y_c = c_1 e^{-2x} + c_2 e^x$. No overlap with $e^{4x}$, so assume $y_p = (Ax + B)e^{4x}$. Computing $y_p', y_p''$ and substituting:
\[
y_p'' + y_p' - 2y_p = 18Ax\,e^{4x} + (9A + 18B)e^{4x} \stackrel{\text{set}}{=} xe^{4x}
\]
gives $18A = 1$, $9A + 18B = 0$, so $A = \tfrac{1}{18}$, $B = -\tfrac{1}{36}$. General solution: $y = (\tfrac{1}{18}x - \tfrac{1}{36})e^{4x} + c_1 e^{-2x} + c_2 e^x$.
 
\textbf{Guess templates:}
\begin{itemize}[leftmargin=*]
    \item $g(x) = $ polynomial of degree $n$: try $A_n x^n + \cdots + A_1 x + A_0$.
    \item $g(x) = e^{\alpha x}$: try $A e^{\alpha x}$.
    \item $g(x) = \sin(\beta x)$ or $\cos(\beta x)$: try $A\sin(\beta x) + B\cos(\beta x)$.
    \item Products and sums: multiply/add the corresponding guesses.
\end{itemize}
 
\textit{i.e.}\ for $y^{(4)} - 5y^{(3)} = 2x^2$: characteristic equation $r^4 - 5r^3 = r^3(r-5) = 0$, so $y_c = c_1 + c_2 x + c_3 x^2 + c_4 e^{5x}$. Since $1, x, x^2$ all appear in $y_c$, multiply the polynomial guess by $x^3$: $y_p = (Ax^2 + Bx + C)x^3$.
 
\section*{Variation of Parameters}
For $y'' + p(x)y' + q(x)y = g(x)$ with $y_c = c_1 y_1 + c_2 y_2$, look for $y_p = u_1(x)y_1 + u_2(x)y_2$ where
\[
u_1' = \frac{-y_2\, g}{W}, \qquad u_2' = \frac{y_1\, g}{W}, \qquad W = W[y_1, y_2].
\]
Note $g$ here is the RHS of the equation in \emph{standard form} (coefficient of $y''$ equal to 1).
 
\textit{i.e.}\ for $y'' + 16y = 3\sec(4x)$: $r^2 + 16 = 0$ gives $r = \pm 4i$, so $y_1 = \sin(4x)$, $y_2 = \cos(4x)$, $W = -4$. Then
\[
u_1' = \frac{-\cos(4x)\cdot 3\sec(4x)}{-4} = \tfrac{3}{4} \;\Rightarrow\; u_1 = \tfrac{3}{4}x,
\]
\[
u_2' = \frac{\sin(4x)\cdot 3\sec(4x)}{-4} = -\tfrac{3}{4}\tan(4x) \;\Rightarrow\; u_2 = \tfrac{3}{16}\ln|\cos(4x)|.
\]
General solution: $y = \tfrac{3}{4}x\sin(4x) + \tfrac{3}{16}\ln|\cos(4x)|\cos(4x) + c_1\sin(4x) + c_2\cos(4x)$.
 
\section*{Euler Equations}
For $x^2 y'' + \alpha x y' + \beta y = 0$ on $(0, \infty)$, try $y = x^r$. This gives the indicial equation $r(r-1) + \alpha r + \beta = 0$.
\begin{itemize}[leftmargin=*]
    \item Distinct real $r_1, r_2$: $y = c_1 x^{r_1} + c_2 x^{r_2}$.
    \item Repeated root $r$: $y = c_1 x^r + c_2 x^r \ln x$.
    \item Complex roots $\alpha \pm \beta i$: $y = x^\alpha\bigl(c_1\cos(\beta\ln x) + c_2\sin(\beta\ln x)\bigr)$.
\end{itemize}
 
\textit{i.e.}\ for $x^2 y'' - 5xy' + 10y = 0$: $r(r-1) - 5r + 10 = r^2 - 6r + 10 = 0$ gives $r = 3 \pm i$, so $y = c_1 x^3 \sin(\ln x) + c_2 x^3 \cos(\ln x)$.
 
\textit{i.e.}\ for $x^2 y'' + 9xy' + 16y = 0$: $r(r-1) + 9r + 16 = r^2 + 8r + 16 = (r+4)^2 = 0$ gives $r = -4$ (repeated), so $y = c_1 x^{-4} + c_2 x^{-4}\ln x$.
 
\section*{Series Solutions at Ordinary Points}
If $x_0$ is an \emph{ordinary point} of $y'' + p(x)y' + q(x)y = 0$ (meaning $p, q$ are analytic at $x_0$), then there exist two linearly independent power-series solutions $y = \sum_{n=0}^\infty a_n (x - x_0)^n$. The radius of convergence is at least the distance from $x_0$ to the nearest singularity of $p$ or $q$ in the complex plane.
 
\textit{i.e.}\ for $(x^2 + 16)y'' + xy' - y = 0$ at $x_0 = 3$: standard form has $\frac{x}{x^2 + 16}$ and $\frac{-1}{x^2 + 16}$ as coefficients, with singularities at $x = \pm 4i$. Distance from 3 to $4i$ is $\sqrt{9 + 16} = 5$, so the radius of convergence is at least 5.
 
\textbf{Procedure.} Substitute $y = \sum a_n x^n$, $y' = \sum n a_n x^{n-1}$, $y'' = \sum n(n-1)a_n x^{n-2}$, reindex so all sums use the same power of $x$, combine into a single sum, and set the coefficient of each $x^n$ to zero to get a recurrence relation.
 
\textit{i.e.}\ for $y'' + 3xy' - y = 0$:
\[
\sum_{n=2}^\infty n(n-1)a_n x^{n-2} + 3\sum_{n=1}^\infty n a_n x^n - \sum_{n=0}^\infty a_n x^n = 0.
\]
Reindexing the first sum and combining: $2a_2 - a_0 + \sum_{n=1}^\infty \bigl[(n+2)(n+1)a_{n+2} + (3n - 1)a_n\bigr]x^n = 0$. So $a_2 = a_0/2$ and
\[
a_{n+2} = \frac{-(3n-1)}{(n+2)(n+1)}\,a_n, \qquad n \geq 1.
\]
This gives a general solution $y = a_0 y_1(x) + a_1 y_2(x)$ with $R = \infty$.
 
\textbf{Shortcut without sigma notation.} If $y(x)$ is unknown and the coefficients of the ODE are themselves functions like $e^x$, expand everything as power series and match coefficients term by term.
 
\textit{i.e.}\ for $y'' + e^x y = 0$ with $y = a_0 + a_1 x + a_2 x^2 + a_3 x^3 + \cdots$ and $y'' = 2a_2 + 6a_3 x + 12 a_4 x^2 + \cdots$:
\[
(2a_2 + 6a_3 x + \cdots) + \bigl(1 + x + \tfrac{x^2}{2!} + \cdots\bigr)(a_0 + a_1 x + a_2 x^2 + \cdots) = 0.
\]
Matching the constant term: $2a_2 + a_0 = 0 \Rightarrow a_2 = -a_0/2$. Matching $x$: $6a_3 + a_1 + a_0 = 0 \Rightarrow a_3 = (-a_1 - a_0)/6$.
 
\section*{Frobenius / Regular Singular Points}
A point $x_0 = 0$ is a \emph{regular singular point} of $y'' + p(x)y' + q(x)y = 0$ if $xp(x)$ and $x^2 q(x)$ are analytic at $0$. Let $p_0 = \lim_{x\to 0} xp(x)$ and $q_0 = \lim_{x\to 0} x^2 q(x)$. Then:
\begin{itemize}[leftmargin=*]
    \item \textbf{Corresponding Euler equation:} $x^2 y'' + p_0 x y' + q_0 y = 0$.
    \item \textbf{Indicial equation:} $r(r-1) + p_0 r + q_0 = 0$. Its roots are the \emph{exponents of the singularity}.
\end{itemize}
Look for solutions of the form $y = \sum_{n=0}^\infty a_n x^{n+r}$.
 
\textit{i.e.}\ for $2x^2 y'' - (3x + x^7)y' + (5\cos^3 x)y = 0$ on $x > 0$: standard form gives $p(x) = -\frac{3x + x^7}{2x^2}$, $q(x) = \frac{5\cos^3 x}{2x^2}$. Then $p_0 = \lim_{x\to 0} x p(x) = -3/2$ and $q_0 = \lim_{x\to 0} x^2 q(x) = 5/2$. The corresponding Euler equation is $x^2 y'' - \tfrac{3}{2}xy' + \tfrac{5}{2}y = 0$, the indicial equation is $r(r-1) - \tfrac{3}{2}r + \tfrac{5}{2} = 0$, i.e.\ $2r^2 - 5r + 5 = 0$, giving exponents $r = \tfrac{5}{4} \pm \tfrac{i\sqrt{15}}{4}$.
 
\textit{i.e.}\ for $4x^2 y'' + (1 + x)y = 0$ on $x > 0$: substituting $y = \sum a_n x^{n+r}$ and matching the lowest power gives the indicial equation $4r^2 - 4r + 1 = 0$, so $r = 1/2$ (double root). The recurrence becomes
\[
a_n = \frac{-a_{n-1}}{(2n+1)(2n-1) + 1}, \qquad n \geq 1.
\]
First terms: $a_1 = -a_0/4$, $a_2 = a_0/(16\cdot 4)$, $a_3 = -a_0/(36\cdot 16 \cdot 4)$, so
\[
y(x) = \sqrt{x}\left(1 - \tfrac{1}{4}x + \tfrac{1}{64}x^2 - \tfrac{1}{2304}x^3 + \cdots\right).
\]
 
\section*{Laplace Transforms --- Definition}
\[
\Lap\{f(t)\} = \int_0^\infty e^{-st} f(t)\,dt
\]
provided the integral converges.
 
\textit{i.e.}\ for $f(t) = 2e^{-7t}$:
\[
\Lap\{2e^{-7t}\} = \int_0^\infty 2e^{-st}e^{-7t}\,dt = 2\int_0^\infty e^{-t(s+7)}\,dt = \frac{2}{-(s+7)}e^{-t(s+7)}\Big|_0^\infty = \frac{2}{s+7}, \quad s > -7.
\]
 
\section*{Common Laplace Transforms}
For any real $a$, positive integer $n$, $c > 0$, and $t_0 \geq 0$:
\begin{align*}
\Lap\{1\} &= \tfrac{1}{s}, \; s > 0 & \Lap\{t^n\} &= \tfrac{n!}{s^{n+1}}, \; s > 0 \\
\Lap\{e^{at}\} &= \tfrac{1}{s-a}, \; s > a & \Lap\{t^n e^{at}\} &= \tfrac{n!}{(s-a)^{n+1}}, \; s > a \\
\Lap\{\sin(at)\} &= \tfrac{a}{s^2 + a^2}, \; s > 0 & \Lap\{\cos(at)\} &= \tfrac{s}{s^2 + a^2}, \; s > 0 \\
\Lap\{u_c(t)\} &= \tfrac{e^{-sc}}{s}, \; s > 0 & \Lap\{\delta(t - t_0)\} &= e^{-st_0}, \; s > 0
\end{align*}
 
\textbf{Translation theorems.} Let $F(s) = \Lap\{f(t)\}$. For $c > 0$:
\begin{itemize}[leftmargin=*]
    \item $\Lap\{e^{ct} f(t)\} = F(s - c)$.
    \item $\Lap\{u_c(t) f(t - c)\} = e^{-cs}F(s)$.
\end{itemize}
 
\textbf{Derivative formulas.} $\Lap\{y'\} = sY - y(0)$, $\Lap\{y''\} = s^2 Y - sy(0) - y'(0)$.
 
\section*{Solving IVPs with Laplace Transforms}
Take the Laplace transform of both sides, solve algebraically for $Y(s)$, decompose into known transforms (often via partial fractions), and invert.
 
\textit{i.e.}\ for $y'' - 2y' - 3y = 6e^{4t}$, $y(0) = -3$, $y'(0) = 1$:
\[
(s^2 Y + 3s - 1) - 2(sY + 3) - 3Y = \tfrac{6}{s - 4} \;\Longrightarrow\; (s^2 - 2s - 3)Y = -3s + 7 + \tfrac{6}{s-4}.
\]
So $Y = \dfrac{-3s^2 + 19s - 22}{(s-3)(s+1)(s-4)}$. Partial fractions: $\dfrac{A}{s-3} + \dfrac{B}{s+1} + \dfrac{C}{s-4}$ with $A = -2$, $B = -11/5$, $C = 6/5$. Inverting:
\[
y(t) = -2e^{3t} - \tfrac{11}{5}e^{-t} + \tfrac{6}{5}e^{4t}.
\]
 
\textit{i.e.}\ to find $\Lap^{-1}\!\left\{\dfrac{s-2}{s^2 + 9} + \dfrac{3}{s^5}\right\}$: split as
\[
\Lap^{-1}\!\left\{\tfrac{s}{s^2 + 9}\right\} - \tfrac{2}{3}\Lap^{-1}\!\left\{\tfrac{3}{s^2 + 9}\right\} + \tfrac{3}{4!}\Lap^{-1}\!\left\{\tfrac{4!}{s^5}\right\} = \cos(3t) - \tfrac{2}{3}\sin(3t) + \tfrac{3}{4!}t^4.
\]
 
\section*{Autonomous Equations and Stability}
An equation $\dydx = f(y)$ is autonomous. The \emph{critical points} (equilibrium solutions) are roots of $f(y) = 0$. To classify, examine the sign of $f$ on either side of each critical point:
\begin{itemize}[leftmargin=*]
    \item Sign goes $+ \to -$: \emph{stable} (asymptotically attracting).
    \item Sign goes $- \to +$: \emph{unstable} (repelling).
    \item Same sign on both sides: \emph{semi-stable}.
\end{itemize}
 
\textit{i.e.}\ for $\dydx = -3y(y^2 - 7y + 10) = -3y(y-5)(y-2)$: critical points are $y = 0, 2, 5$. Testing signs: $y = 0$ and $y = 5$ are stable; $y = 2$ is unstable.
 
\section*{Existence and Uniqueness (Linear First-Order)}
For the linear IVP $\dydx + p(x)y = g(x)$, $y(x_0) = y_0$: the largest interval of existence and uniqueness is the largest interval containing $x_0$ on which both $p$ and $g$ are continuous.
 
\textit{i.e.}\ for $\dydx - \frac{5}{\sin^3 x}y = \frac{6}{x - 2}$, $y(1) = 2$: $-5/\sin^3 x$ is discontinuous at $x = n\pi$ ($n \in \mathbb{Z}$), and $6/(x - 2)$ is discontinuous at $x = 2$. The largest interval containing $x_0 = 1$ on which both are continuous is $(0, 2)$.
 

 
\end{document}
 
